\section{Related Work}
CLIP \cite{radford2021learning} provides a strong English text-image representation. To expand the language of CLIP model,
there are prior studies on learning a bilingual text-image representation ~\cite{ko2022large, bianchi2021contrastive}, and multilingual text-image representation ~\cite{aggarwal2020towards, carlsson2022cross}. In the domain of Chinese text-image pretraining models, prior work such as Taiyi ~\cite{fengshenbang}, CN-CLIP ~\cite{yang2022chinese}, Wukong ~\cite{gu2022wukong}, R2D2 ~\cite{xie2022zero} and BriVL ~\cite{huo2021wenlan, fei2021wenlan}. These methods often need large-scale Chinese text-image pairs and suffer from a significant performance decline in English tasks.
%Most researchers mainly collect large-scale bilingual text-image pairs ($>$100M) to learn the bilingual text encoder from scratch and freeze the parameters of the image encoder.  These works often achieve unsatisfactory performance on zero-shot classifications due to the limited quantity and quality of data beyond English. 

\newcite{carlsson2022cross} proposed a way to utilize Teacher Learning (a.k.a. Knowledge Distillation) \cite{hinton2015distilling} to train a new textual encoder from the original CLIP model with only machine-translated parallel data, without using text-image pairs. Although this method achieves promising cross-lingual retrieval performances with only text data, its zero-shot classification performance in English drops significantly. We follow their work to learn a bilingual model from CLIP with a new text encoder, with the following changes: firstly, we use knowledge distillation on English text pairs in addition to machine-translated text pairs; secondly, we add human-curated translation data for better quality; lastly, we fine-tune the model with text-image pairs to further boost its performance.

%2. Multi-lingual CLIP, Chinese CLIP \\
%In the bilingual setting (Korean and English \cite{ko2022large}, Italian  \cite{bianchi2021contrastive}.

XLM-R \cite{conneau2020unsupervised} is a multilingual language model that achieves strong performances on a wide range of cross-lingual tasks. In our work, we use the XLM-R model as the underlying text encoder and align it with the image encoder trained in CLIP, to achieve competitive performances on cross-lingual and cross-modality tasks. 
%chinese-roberta-wwm \cite{cui2021pre}
